\section{调度案例：从执行轨迹到策略选择}

前面已经说明了常见调度策略的基本规则。要进一步理解这些规则，还需要把任务的到达、
运行、阻塞和完成放在同一条时间线上。调度器观察不到“一个程序总共要运行多久”，它只能
看到已经到达的 task、过去消耗的时间以及当前是否可运行。这一限制决定了理论最优算法与
可实现算法之间的差距。

\subsection{CPU burst 与 I/O burst}

线程从开始到结束通常不会连续占用 CPU。它先执行一段指令，随后等待文件、网络、锁或
定时器；等待条件满足后，再执行下一段指令。两次阻塞之间连续占用处理器的区间称为
\emph{CPU burst}，阻塞等待外部条件的区间称为\emph{I/O burst}。这里的 I/O 是广义说法，
也包括锁等待和定时等待。

\begin{lstlisting}
RUNNABLE -> RUNNING
             |
             +-- CPU burst --+
                              |
                         BLOCKED
                              |
                         I/O burst
                              |
                         RUNNABLE
\end{lstlisting}

这个划分解释了“CPU 密集型”和“I/O 密集型”的准确含义。CPU 密集型任务通常具有较长的
CPU burst，很少主动阻塞；I/O 密集型或交互任务的 CPU burst 较短，但会频繁等待外部事件。
两者都可能运行很长时间，差异在于占用 CPU 与离开 CPU 的节奏。

\begin{longtable}{@{}L{0.18\linewidth}L{0.34\linewidth}L{0.4\linewidth}@{}}
\toprule
工作负载 & 典型执行节奏 & 调度时需要关注的目标 \\
\midrule
批量计算 & 少量长 CPU burst & 吞吐、平均周转时间、cache 局部性 \\
交互请求 & 多次短 CPU burst 与网络等待交替 & 首次响应、p95/p99 延迟、唤醒抢占 \\
后台写回 & 周期性被唤醒，处理一批脏页后再睡眠 & 对前台任务的干扰、I/O 与 CPU 节奏 \\
实时控制 & 周期到达，每次需要有上界的 CPU 时间 & 截止期限、最坏阻塞时间、准入控制 \\
\bottomrule
\end{longtable}

调度器只应在 runnable task 之间分配 CPU。线程处于 I/O burst 时已经离开 runqueue，
它的等待时间由设备、锁持有者或定时器决定。把 blocked task 也计入调度候选，会混淆
“没有处理器时间”和“当前根本不能继续执行”这两种状态。

\subsection{已知服务时间时的理论基准}

先采用一个现实中通常无法满足、但便于建立基准的假设：系统预先知道每个任务还需要多少
CPU 时间。A、B、C 同时在时刻 0 到达，服务时间分别为 8、4、2。若按 A、B、C 的到达
顺序执行，时间线为：

\begin{lstlisting}
time: 0                8        12   14
      |------- A -------|--- B ---| C |
\end{lstlisting}

三个任务的等待时间分别为 0、8、12，平均等待时间为 $20/3$；周转时间分别为 8、12、14，
平均周转时间为 $34/3$。长任务 A 位于队首，使两个短任务都要等待，这就是 convoy effect
在时间指标上的表现。

如果先运行服务时间最短的 C，再运行 B 和 A，时间线变为：

\begin{lstlisting}
time: 0   2        6                14
      | C |--- B ---|------- A -------|
\end{lstlisting}

\begin{longtable}{@{}L{0.16\linewidth}L{0.22\linewidth}L{0.24\linewidth}L{0.28\linewidth}@{}}
\toprule
顺序 & 平均等待时间 & 平均周转时间 & 直接代价 \\
\midrule
A、B、C & $20/3$ & $34/3$ & 短任务排在长任务之后 \\
C、B、A & $8/3$ & $22/3$ & 必须预先知道或准确估计服务时间 \\
\bottomrule
\end{longtable}

在所有任务同时到达、任务之间独立且服务时间已知的条件下，shortest job first（SJF）能
最小化平均等待时间。证明思路可以从交换相邻任务得到。若两个相邻任务 X、Y 的服务时间
满足 $S_X>S_Y$，先运行 X 会使 Y 多等待 $S_X$；交换后，X 只多等待 $S_Y$。由于
$S_Y<S_X$，交换能降低总等待时间。不断消除这样的逆序，最终得到按服务时间递增的顺序。

这个结论有严格的适用前提。它没有保证每个任务都公平，也没有处理任务陆续到达和服务时间
未知的情况。一个持续到来的短任务流可能使长任务长期等待，因此真实系统还要引入 aging、
公平核算或最大延迟约束。

\subsection{任务陆续到达时为何需要抢占}

现在令 A 在时刻 0 到达，需要 8 个单位；B 在时刻 1 到达，需要 4 个单位；C 在时刻 2
到达，需要 2 个单位。若 SJF 不允许抢占，A 一旦开始便运行到时刻 8，B 和 C 仍然要等待。
若每当新任务到达时比较“剩余服务时间”，则得到 shortest remaining time first（SRTF）：

\begin{lstlisting}
time: 0 1 2   4      7              14
      |A|B| C |-- B --|------ A ------|
\end{lstlisting}

\begin{longtable}{@{}L{0.13\linewidth}L{0.18\linewidth}L{0.2\linewidth}L{0.2\linewidth}L{0.2\linewidth}@{}}
\toprule
任务 & 到达 & 完成 & 周转时间 & 等待时间 \\
\midrule
A & 0 & 14 & 14 & $14-8=6$ \\
B & 1 & 7 & 6 & $6-4=2$ \\
C & 2 & 4 & 2 & $2-2=0$ \\
\bottomrule
\end{longtable}

时刻 1，A 还剩 7 个单位，而新到达的 B 只需 4 个，因此 B 抢占 A；时刻 2，B 还剩
3 个单位，而 C 只需 2 个，因此 C 再次抢占。每次决定都基于当前已知任务的剩余时间。

抢占并非没有成本。内核要进入调度路径、保存和恢复现场，工作集也可能离开 cache。若新任务
只比当前任务短很少，抢占节省的等待时间可能小于切换造成的损失。工程实现通常会设置最小
运行粒度、唤醒抢占阈值或延迟容忍度，避免在相近候选者之间频繁切换。

\subsection{未来服务时间为何只能估计}

普通程序不会在每个 CPU burst 开始前声明准确长度。服务时间还会受到输入数据、cache
命中、分支路径、锁竞争和缺页的影响。因此，SJF/SRTF 更适合作为评价基准，实际调度器只能
根据过去行为预测未来。

一种基础预测方法是指数平滑。设第 $n$ 次实际 CPU burst 为 $t_n$，此前预测为 $\tau_n$，
则下一次预测为

\[
\tau_{n+1}=\alpha t_n+(1-\alpha)\tau_n,\qquad 0\leq\alpha\leq 1.
\]

$\alpha$ 越大，最近一次 burst 的影响越强；$\alpha$ 越小，历史平均越稳定。假设初始预测
为 6，$\alpha=0.5$，接下来实际观察到 8 和 2，则预测依次变为

\[
\tau_1=0.5\times 8+0.5\times 6=7,\qquad
\tau_2=0.5\times 2+0.5\times 7=4.5.
\]

预测能识别长期偏交互或偏计算的行为，却无法保证下一次 burst 与过去相似。程序阶段变化
时，旧历史还会造成滞后。MLFQ 采用的“用完整时间片就降低队列级别、很快阻塞则保留较高
级别”，本质上也是一种基于近期行为的在线分类。

\subsection{时间片同时控制响应与开销}

Round-robin 不预测任务长度，而是为每个 runnable task 提供最长为 $Q$ 的连续运行机会。
若有 $N$ 个始终可运行的同权任务，忽略中断和切换成本，某任务再次得到 CPU 的最长队列
等待约为 $(N-1)Q$。减小 $Q$ 能缩短等待上界，但会增加调度频率。

设每次实际任务切换的平均直接成本为 $S$。在任务每次都用满时间片的近似条件下，切换开销
占比约为

\[
\rho=\frac{S}{Q+S}.
\]

若 $S=5\,\mu s$，当 $Q=10\,ms$ 时，$\rho\approx0.05\%$；若把时间片缩短为
$50\,\mu s$，$\rho\approx9.1\%$。后者还没有计入 cache 和 TLB 局部性损失。时间片不能
仅根据“越短响应越快”来选择。

\begin{longtable}{@{}L{0.2\linewidth}L{0.34\linewidth}L{0.38\linewidth}@{}}
\toprule
时间片范围 & 主要收益 & 主要风险 \\
\midrule
相对较长 & 切换较少，计算任务保持局部性 & runnable task 的等待上界增大 \\
相对较短 & 交互任务更快获得第一次服务 & 调度、中断和 cache 恢复成本上升 \\
按任务动态调整 & 可结合延迟目标和任务权重 & 状态与策略更复杂，参数错误会产生抖动 \\
\bottomrule
\end{longtable}

队列边界也必须定义清楚。同一时刻可能同时发生时间片到期、新任务到达和阻塞任务被唤醒。
实现需要规定这些事件的入队次序，否则同一算例可能得到不同结果。多核系统还要决定任务进入
哪个 CPU 的 runqueue；因此，“轮转”不仅是一条 FIFO 规则，还包括计时、事件排序和放置
策略。

\section{公平、优先级与期限的不同承诺}

调度策略经常同时使用“公平”“优先级”“期限”三个词。它们对应不同承诺：公平关心较长
时间窗口内的服务比例，优先级关心当前谁应先获得 CPU，期限关心工作能否在指定时刻前完成。
把三者混为一个数值，会使策略无法解释。

\subsection{加权公平如何换算为服务比例}

设两个始终可运行的任务 A、B 的权重分别为 1024 和 512。若一个 30ms 的观察窗口全部
用于这两个任务，理想服务量分别为

\[
30ms\times\frac{1024}{1024+512}=20ms,\qquad
30ms\times\frac{512}{1024+512}=10ms.
\]

高权重表示在竞争时获得更大比例，并不表示低权重任务永远不能运行。公平调度器需要记录
“已经得到多少加权服务”，以便选择相对欠缺服务的任务。虚拟运行时间可以写成

\[
\Delta v=\Delta t\times\frac{W_0}{w},
\]

其中 $\Delta t$ 是实际运行时间，$w$ 是任务权重，$W_0$ 是基准权重。A 运行 20ms、
B 运行 10ms 后，若 $W_0=1024$，两者的虚拟增量都为 20ms。虚拟时间相同意味着二者
按照各自权重获得了相称服务。

\begin{longtable}{@{}L{0.16\linewidth}L{0.2\linewidth}L{0.22\linewidth}L{0.28\linewidth}@{}}
\toprule
任务 & 权重 & 实际运行 & 虚拟运行增量 \\
\midrule
A & 1024 & 20ms & $20\times1024/1024=20$ \\
B & 512 & 10ms & $10\times1024/512=20$ \\
\bottomrule
\end{longtable}

公平只在存在竞争时分配比例。如果 A 睡眠而 B 独自可运行，让 B 占满 CPU 不违反公平；
空闲处理器没有必要为尚未就绪的 A 保留时间。A 被唤醒后如何放置虚拟时间则需要限制：
若把它放到所有任务之前，频繁睡眠可以获得不受限的抢占优势；若完全忽略睡眠期间的等待，
交互任务又会出现明显延迟。调度器需要在长期份额与唤醒响应之间设置边界。

\subsection{EEVDF 增加了资格与虚拟期限}

只选择最小虚拟运行时间，主要表达“谁得到的加权服务较少”。当任务具有不同延迟需求时，
还需要表达某项工作希望多快得到下一段服务。EEVDF 将决策分为两步：

\begin{enumerate}
\tightlist
\item 根据任务相对于理想服务的 lag，判断它当前是否有资格运行；
\item 在有资格的任务中，选择虚拟截止期最早的任务。
\end{enumerate}

lag 描述任务实际获得的服务与理想份额之间的差额。获得服务不足的任务具有正的补偿需求，
获得过多的任务需要暂时等待。虚拟截止期在公平份额之上表达延迟目标；较短的请求区间会
产生较早的虚拟期限，从而更快得到处理。

资格检查很重要。若只比较 deadline，一个已经超额获得 CPU 的任务可能通过不断提交短请求
长期占先。先限制 eligibility，再比较 deadline，使低延迟选择不能脱离公平核算。这里的
deadline 是调度器虚拟时间中的排序字段，不等同于实时任务必须满足的物理时刻。

\subsection{固定优先级为何会产生反转}

设高优先级任务 H 需要一把 mutex，低优先级任务 L 已经持有该 mutex，中优先级任务 M
与这把锁无关。若 H 到达后阻塞在 L 持有的锁上，而 M 持续抢占 L，就会出现以下关系：

\begin{lstlisting}
L holds mutex
H waits for mutex held by L
M preempts L because M has higher priority than L
\end{lstlisting}

表面上是 M 的优先级低于 H，实际却因为 H 依赖 L，M 延迟了 L，也间接延迟了 H。高优先级
任务的完成时间被低优先级持锁者控制，这称为\emph{优先级反转}。

\begin{longtable}{@{}L{0.18\linewidth}L{0.34\linewidth}L{0.4\linewidth}@{}}
\toprule
阶段 & 未采用继承时 & 采用优先级继承时 \\
\midrule
H 请求 mutex & H 阻塞，等待 L & H 阻塞，并把有效优先级传给 L \\
M 变为 runnable & M 抢占 L & L 以继承后的优先级继续运行 \\
L 释放 mutex & 释放时刻可能被 M 长期推迟 & L 尽快完成临界区并唤醒 H \\
释放之后 & H 才有机会继续 & L 恢复原优先级，H 按高优先级运行 \\
\bottomrule
\end{longtable}

优先级继承不能消除临界区本身的执行时间，但能限制无关中优先级任务造成的额外阻塞。若锁
形成嵌套依赖，继承还要沿等待链传播。实时分析因此必须同时知道调度优先级、锁依赖图和每个
临界区的最坏执行时间。

\subsection{实时期限需要准入与超支处理}

周期任务通常用三项参数描述：周期 $T$、每周期最多需要的运行时间 $C$、相对截止期 $D$。
当 $D=T$ 时，一个任务的处理器利用率为 $U=C/T$。例如三项任务分别为

\[
(C,T)=(2,10),(1,5),(4,20),
\]

总利用率为

\[
U=\frac{2}{10}+\frac{1}{5}+\frac{4}{20}=0.6.
\]

利用率小于 1 只说明总服务需求没有超过单 CPU 的理论容量。共享锁、不可抢占区、中断和
切换成本仍会占用时间；若截止期短于周期，还要使用更完整的可调度性分析。因此，准入控制
必须保留安全余量，并计入系统开销与已知阻塞上界。

任务还可能超过声明的运行预算。若不限制超支，一项异常任务可以破坏所有已准入任务的期限。
系统通常需要在预算耗尽时进行限流、降低调度等级或报告 deadline miss。准入、运行时核算
和超支处理共同构成期限承诺，仅在入队时按 deadline 排序并不足够。

\section{多核调度：处理器位置也是调度结果}

单核调度只选择“运行谁”；多核调度还要选择“在哪里运行”。同一时刻多个 CPU 会独立取
任务、接收唤醒和处理定时器事件。若所有核心共同修改一个全局 runqueue，正确性较直观，
但队列锁会成为频繁竞争的共享 cache line。

\subsection{从全局队列到每 CPU runqueue}

每 CPU runqueue 使多数选择操作只访问本地状态。CPU0 调度时无需取得 CPU1 的队列锁，
不同核心可以并行选择任务。代价是负载可能不均：CPU0 的队列中有多个 runnable task，
CPU1 却可能进入 idle。

\begin{longtable}{@{}L{0.2\linewidth}L{0.34\linewidth}L{0.38\linewidth}@{}}
\toprule
组织方式 & 主要收益 & 必须额外解决的问题 \\
\midrule
单一全局 runqueue & 所有 CPU 看到同一候选集合，天然均衡 & 锁竞争、cache line 迁移、选择结构扩展性 \\
每 CPU runqueue & 本地入队与选择可以并行，保持局部性 & 任务放置、负载均衡、跨 CPU 唤醒与迁移 \\
\bottomrule
\end{longtable}

由此可见，每 CPU 队列没有消除全局调度问题，而是把连续的全局竞争转化为周期性或按需的
协调。系统需要在“少迁移以保持局部性”和“及时迁移以消除空闲”之间取得平衡。

\subsection{新建与唤醒任务放到哪个 CPU}

任务第一次进入 runnable 状态时，内核需要从允许集合中选择目标 CPU。允许集合首先受
affinity、cpuset、安全隔离和 CPU online 状态约束。在剩余候选中，还要考虑队列负载、
处理器容量、cache 拓扑和任务上次运行位置。

被唤醒的任务通常在原 CPU 上保留热 cache、TLB 和 NUMA 内存局部性，因此优先返回原 CPU
可能降低恢复成本。但若原 CPU 正在运行重要任务而另一个 CPU 空闲，保留局部性会增加
runnable wait。唤醒放置需要比较两类成本，而不能只选择队列最短者。

\begin{longtable}{@{}L{0.19\linewidth}L{0.34\linewidth}L{0.39\linewidth}@{}}
\toprule
放置依据 & 倾向保留原 CPU 的原因 & 倾向迁移的原因 \\
\midrule
运行队列 & 避免跨队列加锁和通知 & 原 CPU 的 runnable wait 明显更高 \\
cache/TLB & 上次工作集可能仍在本地 & 空闲 CPU 可立即执行任务 \\
NUMA & 内存页靠近原 CPU 所在节点 & 任务已经迁移内存或远端等待过长 \\
CPU capacity & 原核心可能性能更适合当前任务 & 高容量核心被更紧急任务占用 \\
能耗 & 避免唤醒额外核心 & 集中负载可能触发降频或热限制 \\
\bottomrule
\end{longtable}

异构处理器进一步使“一个 runnable task 等于一份相同负载”的假设失效。高性能核心和节能
核心单位时间能完成的工作不同，任务还可能受到热设计功耗和频率上限影响。负载比较需要按
CPU capacity 归一化，并结合任务利用率估计。

\subsection{迁移必须是原子的队列转移}

任务从 CPU0 迁移到 CPU1 时，不能先在目标队列插入、稍后再从源队列删除，否则两个 CPU
可能同时选择同一个 task；也不能先解除所有归属、再无保护地插入，否则并发唤醒可能使
任务丢失。迁移过程至少要保持以下顺序：

\begin{enumerate}
\tightlist
\item 按固定次序取得源 runqueue 与目标 runqueue 的锁；
\item 确认 task 仍位于源队列，状态仍为 runnable，且没有正在某个 CPU 上执行；
\item 从源队列删除 task，更新 CPU 归属和迁移统计；
\item 将 task 插入目标队列，必要时通知目标 CPU 重新调度；
\item 释放两端队列锁。
\end{enumerate}

\begin{longtable}{@{}L{0.22\linewidth}L{0.32\linewidth}L{0.38\linewidth}@{}}
\toprule
若边界处理错误 & 可能出现的状态 & 外部表现 \\
\midrule
先入目标、后出源 & task 同时位于两个 runqueue & 同一内核栈被两个 CPU 恢复，状态遭到破坏 \\
先出源、丢失引用 & runnable task 不在任何队列 & 线程永久得不到 CPU \\
双队列锁顺序不一致 & CPU0 等 CPU1 的锁，CPU1 等 CPU0 的锁 & 调度路径死锁 \\
迁移 running task & task 仍在源 CPU 执行又被目标选择 & 两个 current 指向同一执行现场 \\
\bottomrule
\end{longtable}

这里的核心不变量仍然是“一个 runnable task 只属于一个 runqueue”，只是多核环境允许多个
CPU 同时修改相关状态，因此状态检查、队列转移和归属更新必须成为一个受保护的整体。

\subsection{主动均衡与空闲拉取}

负载均衡可以由繁忙 CPU 主动把任务推走，也可以由即将 idle 的 CPU 从其他队列拉取任务。
主动推送能较早缓解长队列，但需要扫描其他 CPU 并承担跨核通知；空闲拉取只在确有空闲容量
时发生，却可能让任务先经历一段不必要的等待。

\begin{longtable}{@{}L{0.2\linewidth}L{0.34\linewidth}L{0.38\linewidth}@{}}
\toprule
触发方式 & 适用时机 & 需要避免的问题 \\
\midrule
周期均衡 & 定期检查相邻调度域的负载差 & 周期过短导致迁移抖动，过长导致空闲与排队并存 \\
idle balance & CPU 即将无任务可运行 & 扫描成本超过即将获得的运行收益 \\
wakeup balance & blocked task 重新变为 runnable & 每次唤醒都跨核迁移，破坏 cache 局部性 \\
overload push & 实时或 deadline 队列出现过载 & 把不可迁移或亲和性受限任务错误移出 \\
\bottomrule
\end{longtable}

调度域按硬件拓扑分层。系统通常先在共享核心或共享 LLC 的近邻间均衡，再考虑同一 NUMA
节点，最后才跨节点或跨 socket。距离越远，可利用的空闲 CPU 越多，迁移带来的 cache 与
内存局部性代价也越高。

\subsection{SMT 线程并不等同于完整核心}

同时多线程（SMT）的两个逻辑 CPU 共享部分执行资源。若两个计算密集任务同时放在同一物理
核心的 sibling 上，它们可能争用执行端口和 cache；把其中一个迁到独立物理核心通常能获得
更高吞吐。另一方面，当系统只剩一个空闲 SMT sibling 时，让任务运行仍可能优于继续排队。

调度器因此需要理解拓扑层级，不能仅按逻辑 CPU 数量平均任务。安全隔离还可能要求互不信任
的工作负载不共享 SMT 核心，这会进一步缩小允许放置集合。性能、公平和隔离在此处共同影响
任务位置。

\subsection{cgroup 配额为何会产生周期性延迟}

设一个服务组的配额为每 100ms 可使用 20ms CPU。组内四个线程在周期开始时同时变为
runnable，并在 20ms 内耗尽共享预算。即使宿主机随后还有空闲 CPU，这四个线程仍可能被
throttle 到下一个周期：

\begin{lstlisting}
0ms             20ms                              100ms
| group runs     |----------- throttled -----------|
                                                    refill
\end{lstlisting}

这项限制针对整个 group 的累计运行时间，不是给每个线程各 20ms。若请求在 21ms 到达，
它可能等待接近 79ms 才能继续，因此形成接近 quota period 的尾延迟尖峰。缩短 period
可以降低最长等待，却增加配额核算和补充频率；增加 quota 则会改变租户之间的资源隔离。

层级 cgroup 还要求同时满足祖先和当前组的预算。子组仍有余额时，若父组预算已经耗尽，
子组也不能运行。解释容器中的 CPU 延迟时，需要沿层级检查配额、实际消耗和 throttled
时间，单看线程优先级无法得到完整原因。

\section{进程与线程生命周期的完整边界}

调度器选择的是 task，但 task 引用的进程资源有各自生命周期。创建、exec、线程退出和
进程回收会改变这些引用关系。若只维护调度状态而忽略资源可见性，新 task 可能访问尚未
初始化的对象，退出 task 也可能留下无法回收或过早释放的资源。

\subsection{创建 task 时先构造，后发布}

创建新 task 可以分为私有构造阶段和公开发布阶段。构造阶段分配 task 与内核栈，准备用户
返回现场，复制或共享地址空间、文件表、凭据和信号状态；这些步骤尚未完成时，新对象只能由
创建路径访问。发布阶段分配可见标识、连接父子关系，并把 task 置为 runnable。

\begin{longtable}{@{}L{0.18\linewidth}L{0.36\linewidth}L{0.38\linewidth}@{}}
\toprule
阶段 & 可见范围 & 失败时的处理 \\
\midrule
分配 task 与栈 & 仅创建路径持有局部引用 & 释放已经分配的内存 \\
建立资源引用 & 尚未进入 PID 查找和 runqueue & 按逆序撤销 mm、files、cred 等引用 \\
建立进程关系 & 父进程与管理结构开始可见 & 在锁保护下解除关系和标识 \\
进入 runqueue & 任意 CPU 可能立即选择该 task & 此后不能再依赖“子任务尚未运行”的假设 \\
\bottomrule
\end{longtable}

“最后入队”是发布边界。若 task 提前进入 runqueue，另一个 CPU 可能在文件表、地址空间或
返回寄存器尚未准备好时执行它。发布之后再补字段不能保证安全，因为调度器没有义务等待
创建者完成剩余工作。

\subsection{共享选择决定线程与进程的边界}

创建接口可以针对不同资源选择“复制”或“共享”。共享地址空间使两个 task 成为同一进程
中的线程；共享文件表会使一方关闭 fd 直接影响另一方；共享信号处理表则使处理动作属于
共同资源。线程边界由一组共享关系共同形成，不能仅由一个名称标志决定。

\begin{longtable}{@{}L{0.2\linewidth}L{0.34\linewidth}L{0.38\linewidth}@{}}
\toprule
资源 & 复制后的语义 & 共享后的语义 \\
\midrule
地址空间 & 修改私有映射通常不影响对方 & 任意线程的普通内存写对其他线程可见 \\
文件描述符表 & 关闭本方 fd 槽不删除对方槽位 & 关闭共享槽位会影响所有引用该表的线程 \\
文件对象 & 不同 fd 表仍可引用同一打开文件对象 & offset、状态标志按对象规则共同变化 \\
信号处理配置 & 后续修改各自独立 & 处理函数配置属于线程组共同状态 \\
凭据 & 后续身份变化按各自对象处理 & 必须满足系统对线程组身份一致性的约束 \\
\bottomrule
\end{longtable}

即使共享地址空间，每个线程仍需要独立用户栈和内核栈。若两个线程使用同一栈，它们的函数
调用、局部变量和返回地址会互相覆盖；若共享一套寄存器现场，调度器也无法分别暂停和恢复
它们。因此，“共享进程资源”和“保留独立执行现场”必须同时成立。

\subsection{多线程进程中的 exec}

exec 的语义是用新程序替换当前进程的用户地址空间。单线程时，只需准备新映射和初始用户
现场，再在提交点替换旧地址空间。多线程时，其他线程可能正在旧地址空间中执行；若它们在
替换之后继续运行，保存的 RIP、RSP 和用户指针将指向已经失效的映射。

因此，多线程 exec 需要先使调用线程成为唯一能够继续的执行流，等待或终止同一线程组中的
其他线程，再提交新地址空间。提交点之前失败，旧程序应尽量继续运行并收到错误；提交点之后，
旧地址空间已经不可恢复，控制流只能进入新程序或终止。

\begin{longtable}{@{}L{0.18\linewidth}L{0.36\linewidth}L{0.38\linewidth}@{}}
\toprule
阶段 & 仍可返回旧程序吗 & 需要保持的条件 \\
\midrule
读取并验证参数 & 可以 & 旧 mm、旧用户栈和调用线程仍有效 \\
打开并检查可执行文件 & 可以 & 尚未改变线程组的公开执行状态 \\
构造新 mm 与用户栈 & 通常可以 & 新对象尚未成为当前地址空间 \\
处理其他线程 & 进入敏感过渡阶段 & 不能再允许其他线程恢复旧用户现场 \\
提交新 mm & 不再返回旧程序 & 当前 RIP、RSP、凭据和 fd 规则与新程序一致 \\
\bottomrule
\end{longtable}

exec 保留 PID 并不表示所有进程状态都保留。地址空间和用户栈被替换，设置 close-on-exec
的 fd 被关闭，信号处理配置按接口规则重置，凭据可能根据可执行文件发生受控转换。每类
资源都有独立规则，不能用“进程没变”或“进程全换了”概括。

\subsection{线程退出与进程退出}

一个线程退出时，它自己的用户栈映射、内核栈、寄存器现场和调度实体可以进入释放流程，
但同一进程中的其他线程仍可能引用共享地址空间和文件表。只有最后一个相关引用消失时，
共享对象才可以释放。

\begin{lstlisting}
thread exit:
  stop scheduling this task
  release per-thread state
  drop references to shared process objects
  if other threads remain:
    keep mm/files/process identity alive
  else:
    begin process-wide exit and parent notification
\end{lstlisting}

进程级退出还要关闭文件引用、解除地址空间、记录退出码并通知父进程。这里存在两个不同的
完成点：task 已经不再运行，以及父进程已经读取退出状态并完成回收。前者保证执行停止，
后者决定最后的进程记录何时释放。

\subsection{zombie 保存的是尚未交付的退出结果}

子进程停止执行后，父进程仍可能需要退出码、终止原因和资源使用统计。若内核立即删除全部
记录，\code{wait} 将无法区分“子进程尚未退出”“已经退出但结果丢失”和“从来不是自己的
子进程”。因此系统保留一个不再可调度的 zombie 记录。

\begin{longtable}{@{}L{0.22\linewidth}L{0.33\linewidth}L{0.37\linewidth}@{}}
\toprule
状态 & 仍然保留的内容 & 已经释放或禁止的内容 \\
\midrule
RUNNING/RUNNABLE & 完整执行现场和资源引用 & 无 \\
EXITING & 退出码、清理进度、必要引用 & 不再接受新的普通工作 \\
ZOMBIE & PID、父子关系、退出码、必要统计 & 用户地址空间、普通执行资源；不能再次调度 \\
REAPED & 不再保留可查询的子进程结果 & 最后进程记录与 PID 可进入复用流程 \\
\bottomrule
\end{longtable}

父进程调用 wait 后，退出结果完成一次性交付，zombie 才转为 reaped。若父进程提前退出，
系统必须把仍需回收的子进程转交给指定的接管者，确保退出结果最终有人消费。zombie 的存在
是父子进程协议的一部分，不是仍在运行的“空进程”。

\subsection{PID、task 引用与编号复用}

PID 是查找进程的标识，不等同于对象自身。进程被回收后，有限编号空间允许复用 PID。
若内核或管理程序只保存一个整数，稍后再按该整数查找，可能得到新一代进程。需要跨时间稳定
引用时，还应保存对象引用、创建代次或专用进程句柄。

这一问题与设备 request tag 的复用相似：编号只在特定生命周期内唯一。对象进入回收前，
所有能够到达它的异步路径都必须完成或放弃引用；编号可以复用之后，迟到事件不能再作用于
新对象。

\section{本章不变量与综合推演}

至此，可以用一个包含创建、调度、阻塞、迁移和退出的场景检查本章全部边界。父进程 P 在
CPU0 创建子进程 C；C 第一次进入 runqueue 后在 CPU1 运行，读取 socket 时阻塞；数据到达
后 C 被唤醒到 CPU0，运行一个时间片后迁移回 CPU1，最终退出并由 P 回收。

\begin{longtable}{@{}L{0.13\linewidth}L{0.35\linewidth}L{0.44\linewidth}@{}}
\toprule
阶段 & 主要状态变化 & 必须保持的不变量 \\
\midrule
创建 & C 从私有构造对象变为 RUNNABLE & 所有必要资源已初始化，C 只进入一个 runqueue \\
首次调度 & C 从 RUNNABLE 变为 CPU1 的 current & C 已从 runqueue 移除，不在其他 CPU 执行 \\
阻塞读取 & C 从 RUNNING 变为 BLOCKED & 等待谓词与 wait queue 已登记，C 不再占 CPU \\
网络唤醒 & C 从 BLOCKED 变为 RUNNABLE & 唤醒原因已记录，只进入一个目标 runqueue \\
抢占 & C 的现场保存后重新入队 & 状态与 on\_rq 一致，共享进程资源引用不变 \\
迁移 & C 从 CPU0 队列原子转移到 CPU1 队列 & 迁移期间不会同时属于两队列，也不会丢失 \\
退出 & C 停止调度并形成退出结果 & 地址空间与文件引用按引用计数释放，退出码仍可读取 \\
回收 & P 取得退出结果，C 进入 REAPED & 最后引用消失后才释放记录并允许 PID 复用 \\
\bottomrule
\end{longtable}

本章的基础验收可以归纳为以下各项：

\begin{itemize}
\tightlist
\item 进程资源边界与线程执行现场能够分别说明，不能把程序文件、进程和 task 混为一体。
\item RUNNABLE task 必须位于且只位于一个 runqueue，BLOCKED 和 ZOMBIE task 不得被选择。
\item 阻塞前完成谓词检查和等待登记，唤醒只授予运行资格，恢复后仍要重新检查条件。
\item 抢占、系统调用入口和上下文切换分别记录，不能用其中一个事件代替另一个事件的证据。
\item 多核迁移保持源队列、目标队列和 task 归属的一致，不允许重复执行或丢失 runnable task。
\item 公平份额、固定优先级和实时期限按各自承诺评价，锁阻塞与运行预算同时纳入分析。
\item 新 task 完整初始化后才能发布；线程退出与进程回收按共享引用和父子协议分阶段完成。
\end{itemize}

这些不变量不依赖某一种具体调度算法。FCFS、round-robin、MLFQ、公平调度或 deadline
调度可以改变候选 task 的排序，却不能改变状态、队列、执行现场和资源生命周期的基本关系。
